Benefiting from the scaling of training data and key architectural enhancements, Qwen3-VL demonstrates substantially improved video understanding capabilities. In particular, the integration of interleaved MRoPE, the insertion of textual timestamps, and scaling temporally dense video captions collectively enable the Qwen3-VL 8B variant to achieve performance competitive with the significantly larger Qwen2.5-VL 72B model.

We conduct a comprehensive evaluation across a diverse set of video understanding tasks, encompassing general video understanding (VideoMME~\citep{fu2024video}, MVBench~\citep{li2024mvbench}), temporal video grounding (Charades-STA~\citep{gao2017tall}), video reasoning (VideoMMMU~\citep{hu2025video}, MMVU~\citep{zhao2025mmvu}), and long-form video understanding (LVBench~\citep{wang2024lvbench}, MLVU~\citep{zhou2024mlvu}). In comparison with state-of-the-art proprietary models — including Gemini 2.5 Pro, GPT-5, and Claude Opus 4.1, Qwen3-VL demonstrates competitive and, in several cases, superior performance.
In particular, our flagship model, Qwen3-VL-235B-A22B-Instruct, achieves performance on par with leading models such as Gemini 2.5 Pro (with a thinking budget of 128) and GPT-5 minimal on standard video understanding benchmarks. By extending the context window to 256K tokens, it further attains or even surpasses Gemini-2.5-Pro on long-video evaluation tasks, most notably on MLVU. 

Regarding evaluation details, we imposed a cap of 2,048 frames per video for all benchmarks, ensuring that the total number of video tokens did not exceed 224K. The maximum number of tokens per frame was set to 768 for VideoMMMU and MMVU, and to 640 for all other benchmarks. Additionally, videos from Charades-STA were sampled at 4 frames per second (fps), while a rate of 2 fps was used for all other benchmarks. For VideoMMMU, we employed a model-based judge for evaluation, as rule-based scoring proved insufficiently accurate. It is worth noting that our comparison cannot guarantee full fairness due to resource and API limitations, which constrained the number of input frames used during evaluation: 512 for Gemini 2.5 Pro, 256 for GPT-5, and 100 for Claude Opus 4.1. 
